\documentclass[12pt,a4paper]{article}

\usepackage[utf8]{inputenc}
\usepackage[T1]{fontenc}
\usepackage{amsmath, amssymb, amsthm}
\usepackage{graphicx}
\usepackage{hyperref}
\usepackage[letterpaper, margin=2.5cm]{geometry}
\usepackage{array}

\setlength{\parindent}{0pt}
\setlength{\parskip}{6pt}

\hypersetup{
    colorlinks=true,
    linkcolor=blue,
    citecolor=blue,
    urlcolor=blue
}

\begin{document}

\begin{titlepage}
\begin{center}

\vspace*{2cm}

{\LARGE \textbf{Universidad de Chile}}\\[0.8cm]
{\Large Facultad de Ingeniería}\\[0.5cm]
{\Large \textbf{MATE1189 --- Cálculo Avanzado}}\\[2cm]

\rule{\linewidth}{0.5mm}\\[0.8cm]
{\Huge \textbf{Optimización de Paneles Solares}}\\[0.4cm]
{\large Funciones en Varias Variables --- Proyecto Final}\\[0.4cm]
\rule{\linewidth}{0.5mm}\\[2.5cm]

\begin{tabular}{rl}
\textbf{Profesor:} & Dr. Andy Alexis Vidal Yunge \\[0.3cm]
\textbf{Integrantes:} & Joaquín Cantero \\[0.15cm]
& Rodrigo Altamirano \\[0.15cm]
& Diego Álvarez \\[0.3cm]
\textbf{Fecha:} & 14 de junio de 2026
\end{tabular}

\vfill

{\small \textit{Versión preliminar --- sujeta a modificaciones}}

\end{center}
\end{titlepage}

\clearpage

\tableofcontents
\clearpage

%% ====================================================================
\section{Introducción}
%% ====================================================================

\subsection{Contexto}

La energía solar fotovoltaica se ha consolidado como una de las principales fuentes de energía renovable a nivel mundial. La eficiencia de un sistema fotovoltaico depende críticamente de la orientación e inclinación de los paneles solares respecto de la posición del Sol. Una instalación inadecuada puede reducir significativamente la energía captada, afectando la rentabilidad y eficiencia del sistema.

Mediante funciones de varias variables es posible modelar la energía captada por un panel solar en función de parámetros como el ángulo de inclinación y la orientación respecto del norte geográfico. Las herramientas del cálculo multivariable permiten analizar estas relaciones, identificar configuraciones óptimas y estudiar la sensibilidad del sistema ante pequeñas variaciones en los parámetros de instalación.

\subsection{Problema de investigación}

El presente proyecto busca responder la siguiente pregunta:

\begin{quote}
\itshape
¿Cómo puede determinarse la orientación e inclinación óptima de un panel solar utilizando herramientas de cálculo multivariable para maximizar la energía captada por el sistema?
\end{quote}

\subsection{Objetivos}

\subsubsection{Objetivo general}

Diseñar y analizar un modelo matemático que permita optimizar la orientación de paneles solares mediante herramientas de cálculo multivariable.

\subsubsection{Objetivos específicos}

\begin{enumerate}
    \item Construir un modelo matemático para la energía captada por un panel solar.
    \item Analizar la influencia de la inclinación y orientación sobre la energía recibida.
    \item Determinar e interpretar las derivadas parciales del modelo.
    \item Construir e interpretar el gradiente de la función energía.
    \item Identificar configuraciones óptimas mediante técnicas de optimización.
    \item Estudiar curvas de nivel y superficies asociadas al modelo.
    \item Aplicar aproximaciones lineales y derivadas direccionales.
    \item Implementar una simulación computacional que permita analizar distintas configuraciones.
    \item Evaluar la sensibilidad del sistema frente a errores de instalación.
\end{enumerate}

%% ====================================================================
\section{Marco teórico}
%% ====================================================================

\subsection{Funciones de varias variables}

Una función de dos variables $f: D \subseteq \mathbb{R}^2 \to \mathbb{R}$ asigna a cada par $(x,y) \in D$ un único valor real $f(x,y)$. En este proyecto, la energía captada $E$ depende de dos variables geométricas: el ángulo de inclinación $\theta$ y el ángulo de orientación $\phi$.

\subsection{Derivadas parciales}

La derivada parcial de $f$ respecto de $x$ en $(a,b)$ se define como

\[
\frac{\partial f}{\partial x}(a,b) = \lim_{h\to 0} \frac{f(a+h,b)-f(a,b)}{h},
\]

y representa la tasa de variación de $f$ cuando solo $x$ cambia, manteniendo $y$ fijo. Análogamente para $\partial f/\partial y$. Geométricamente, corresponde a la pendiente de la recta tangente a la superficie $z=f(x,y)$ en la dirección del eje respectivo.

\subsection{Gradiente}

El gradiente de $f$ es el vector

\[
\nabla f(x,y) = \left( \frac{\partial f}{\partial x},\; \frac{\partial f}{\partial y} \right).
\]

El gradiente apunta en la dirección de máximo crecimiento de la función y su magnitud indica la tasa de cambio en esa dirección. Si $\nabla f = \mathbf{0}$, el punto es crítico y candidato a extremo local.

\subsection{Derivada direccional}

La derivada direccional en la dirección del vector unitario $\mathbf{u}=(u_1,u_2)$ es

\[
D_{\mathbf{u}} f(x,y) = \nabla f(x,y)\cdot\mathbf{u}
= \frac{\partial f}{\partial x}\,u_1 + \frac{\partial f}{\partial y}\,u_2,
\]

y mide la tasa de cambio de $f$ al desplazarse en la dirección $\mathbf{u}$.

\subsection{Plano tangente y linealización}

La linealización de $f$ en torno a $(a,b)$ es

\[
L(x,y) = f(a,b) + \frac{\partial f}{\partial x}(a,b)\,(x-a) + \frac{\partial f}{\partial y}(a,b)\,(y-b),
\]

cuyo gráfico es el plano tangente a $z=f(x,y)$ en $(a,b,f(a,b))$. Es útil para estimar valores de $f$ cerca de configuraciones conocidas.

\subsection{Puntos críticos y matriz Hessiana}

Un punto $(a,b)$ es crítico si $\nabla f(a,b)=\mathbf{0}$. Para clasificarlo se usa la matriz Hessiana

\[
H_f(a,b) =
\begin{pmatrix}
\frac{\partial^2 f}{\partial x^2} & \frac{\partial^2 f}{\partial x\partial y} \\[6pt]
\frac{\partial^2 f}{\partial y\partial x} & \frac{\partial^2 f}{\partial y^2}
\end{pmatrix}_{(a,b)}.
\]

El criterio de la segunda derivada establece: si $\nabla f(a,b)=\mathbf{0}$ y $D = \det H_f(a,b)$,
\begin{itemize}
    \item $D>0$ y $\partial^2 f/\partial x^2 > 0$: mínimo local;
    \item $D>0$ y $\partial^2 f/\partial x^2 < 0$: máximo local;
    \item $D<0$: punto de silla;
    \item $D=0$: criterio no concluyente.
\end{itemize}

\subsection{Curvas de nivel y superficies}

Las curvas de nivel son $\{(x,y)\in D \mid f(x,y)=c\}$. En el proyecto representan configuraciones $(\theta,\phi)$ con igual rendimiento energético. El gráfico $z=f(x,y)$ es una superficie en $\mathbb{R}^3$.

%% ====================================================================
\section{Modelo matemático}
%% ====================================================================

\subsection{Variables geométricas}

La orientación del panel se describe mediante dos variables angulares:
\begin{itemize}
    \item $\theta$: ángulo de inclinación respecto de la horizontal. Rango físico $\theta\in[0,\pi/2]$ ($\theta=0$: horizontal, $\theta=\pi/2$: vertical).
    \item $\phi$: ángulo de orientación (azimut) respecto del norte geográfico, $\phi\in[0,2\pi)$ ($\phi=0$: panel mirando al norte).
\end{itemize}

\subsection{Modelo simplificado}

Como aproximación inicial se considera la función

\begin{equation}\label{eq:base}
E(\theta,\phi)=A\cos(\theta-\theta_0)\cos(\phi-\phi_0),
\end{equation}

donde:
\begin{itemize}
    \item $A>0$: máxima energía posible de captar.
    \item $\theta_0$: ángulo de inclinación ideal para la ubicación geográfica estudiada.
    \item $\phi_0$: orientación ideal para la ubicación geográfica estudiada.
\end{itemize}

La función modela la disminución de energía cuando el panel se aleja de la orientación óptima: producto de cosenos que alcanza su máximo en $(\theta_0,\phi_0)$ y decae suavemente en todas las direcciones.

\subsection{Modelo extendido}

El modelo simplificado puede enriquecerse incorporando factores adicionales que afectan la radiación recibida.

\subsubsection{Variación diaria de la posición solar}

La posición del Sol cambia durante el día. Esto puede incorporarse haciendo que los ángulos óptimos dependan del tiempo:

\begin{equation}\label{eq:ext_tiempo}
E(\theta,\phi,t)=A\cos\!\big(\theta-\theta_0(t)\big)\cos\!\big(\phi-\phi_0(t)\big),
\end{equation}

donde $\theta_0(t)$ y $\phi_0(t)$ describen la elevación y azimut solar en cada instante.

\subsubsection{Estaciones del año}

La inclinación ideal $\theta_0$ varía a lo largo del año por la declinación solar. Una aproximación es

\begin{equation}\label{eq:ext_estaciones}
\theta_0(d)=\theta_{\text{lat}} + \delta\,\sin\!\left(\frac{2\pi d}{365}\right),
\end{equation}

con $d$ el día del año, $\theta_{\text{lat}}$ determinado por la latitud y $\delta\approx 23.45^\circ$ la declinación máxima.

\subsubsection{Nubosidad}

La nubosidad reduce la radiación incidente mediante un factor de atenuación $N(t)\in[0,1]$:

\begin{equation}\label{eq:ext_nubosidad}
E(\theta,\phi,t)=N(t)\,A\cos(\theta-\theta_0)\cos(\phi-\phi_0).
\end{equation}

\subsubsection{Sombras producidas por obstáculos cercanos}

Edificaciones u obstáculos proyectan sombras que reducen la energía. Se introduce una función de penalización $S(\theta,\phi)\in[0,1]$:

\begin{equation}\label{eq:ext_sombras}
E(\theta,\phi)=A\cos(\theta-\theta_0)\cos(\phi-\phi_0)\,\big[1-S(\theta,\phi)\big].
\end{equation}

\subsubsection{Ubicación geográfica}

La latitud y longitud determinan la trayectoria solar y, por tanto, los valores de $\theta_0$ y $\phi_0$. Un modelo dependiente de la ubicación requiere expresiones basadas en la geometría solar que relacionen $\theta_0$ y $\phi_0$ con la latitud, el día del año y la hora del día.

Combinando estos factores, el modelo extendido tiene la forma general

\begin{equation}\label{eq:ext_integrado}
E(\theta,\phi,t)=A\;N(t)\;\big[1-S(\theta,\phi)\big]\;
\cos\!\big(\theta-\theta_0(t)\big)\cos\!\big(\phi-\phi_0(t)\big),
\end{equation}

donde $t$ representa tanto la hora del día como el día del año.

%% ====================================================================
\section{Desarrollo analítico}
%% ====================================================================

Aplicamos las herramientas del cálculo multivariable al modelo simplificado

\[
E(\theta,\phi)=A\cos(\theta-\theta_0)\cos(\phi-\phi_0),
\]

con $(\theta,\phi)\in[0,\pi/2]\times[0,2\pi)$ y parámetros $A>0$, $\theta_0\in[0,\pi/2]$, $\phi_0\in[0,2\pi)$.

\subsection{Derivadas parciales}

Calculamos las derivadas parciales de primer orden:

\begin{align}
\frac{\partial E}{\partial\theta}(\theta,\phi)
&= -A\,\sin(\theta-\theta_0)\cos(\phi-\phi_0), \label{eq:dEdt}\\[4pt]
\frac{\partial E}{\partial\phi}(\theta,\phi)
&= -A\,\cos(\theta-\theta_0)\sin(\phi-\phi_0). \label{eq:dEdp}
\end{align}

\noindent\textit{Interpretación.} $\partial E/\partial\theta$ mide la sensibilidad de la energía ante pequeños cambios en la inclinación, manteniendo fija la orientación. Análogamente, $\partial E/\partial\phi$ mide la sensibilidad ante cambios en el azimut. Cerca del óptimo estas derivadas son pequeñas; lejos de él, crecen en magnitud.

\subsection{Gradiente}

El gradiente de la función energía es

\begin{equation}\label{eq:grad}
\nabla E(\theta,\phi)=\big(
-A\sin(\theta-\theta_0)\cos(\phi-\phi_0),\;
-A\cos(\theta-\theta_0)\sin(\phi-\phi_0)
\big).
\end{equation}

Su dirección indica hacia dónde debe modificarse la orientación para aumentar la energía lo más rápidamente posible. Su magnitud es la tasa de cambio en esa dirección.

\subsection{Puntos críticos}

Los puntos críticos satisfacen $\nabla E(\theta,\phi)=\mathbf{0}$:

\begin{align}
\sin(\theta-\theta_0)\cos(\phi-\phi_0) &= 0, \label{eq:cp1}\\
\cos(\theta-\theta_0)\sin(\phi-\phi_0) &= 0. \label{eq:cp2}
\end{align}

De \eqref{eq:cp1} se tiene $\sin(\theta-\theta_0)=0$ o $\cos(\phi-\phi_0)=0$.
De \eqref{eq:cp2} se tiene $\cos(\theta-\theta_0)=0$ o $\sin(\phi-\phi_0)=0$.
Analizando las combinaciones:

\begin{itemize}
    \item \textbf{Caso 1:} $\sin(\theta-\theta_0)=0$ y $\sin(\phi-\phi_0)=0$ \\
    $\Rightarrow$ $\theta-\theta_0=n\pi$, $\phi-\phi_0=m\pi$, con $n,m\in\mathbb{Z}$.

    \item \textbf{Caso 2:} $\cos(\phi-\phi_0)=0$ y $\cos(\theta-\theta_0)=0$ \\
    $\Rightarrow$ $\phi-\phi_0=\pi/2+m\pi$, $\theta-\theta_0=\pi/2+n\pi$, $n,m\in\mathbb{Z}$.
\end{itemize}

En el dominio físico los puntos críticos de interés son

\begin{equation}\label{eq:cp_final}
(\theta,\phi)=(\theta_0+n\pi,\;\phi_0+m\pi),\qquad n,m\in\mathbb{Z}.
\end{equation}

\subsection{Matriz Hessiana y clasificación}

Las derivadas parciales de segundo orden son:

\begin{align}
\frac{\partial^2 E}{\partial\theta^2}
&= -A\cos(\theta-\theta_0)\cos(\phi-\phi_0), \\[4pt]
\frac{\partial^2 E}{\partial\theta\partial\phi}
&= A\sin(\theta-\theta_0)\sin(\phi-\phi_0), \\[4pt]
\frac{\partial^2 E}{\partial\phi\partial\theta}
&= A\sin(\theta-\theta_0)\sin(\phi-\phi_0), \\[4pt]
\frac{\partial^2 E}{\partial\phi^2}
&= -A\cos(\theta-\theta_0)\cos(\phi-\phi_0).
\end{align}

La matriz Hessiana es

\begin{equation}\label{eq:hess}
H_E(\theta,\phi)=
\begin{pmatrix}
-A\cos(\theta-\theta_0)\cos(\phi-\phi_0) &
A\sin(\theta-\theta_0)\sin(\phi-\phi_0) \\[6pt]
A\sin(\theta-\theta_0)\sin(\phi-\phi_0) &
-A\cos(\theta-\theta_0)\cos(\phi-\phi_0)
\end{pmatrix}.
\end{equation}

Evaluamos en los puntos críticos $(\theta_0+n\pi,\phi_0+m\pi)$.

\subsubsection{Óptimo $(\theta_0,\phi_0)$}

\[
H_E(\theta_0,\phi_0)=\begin{pmatrix}-A & 0\\0 & -A\end{pmatrix},
\quad
D = A^2 > 0,\quad
\frac{\partial^2 E}{\partial\theta^2} = -A < 0.
\]

Por el criterio de la segunda derivada, $(\theta_0,\phi_0)$ es un \textbf{máximo local}. Físicamente es la configuración que maximiza la energía captada.

\subsubsection{Mínimo $(\theta_0+\pi,\phi_0+\pi)$}

\[
H_E(\theta_0+\pi,\phi_0+\pi)=\begin{pmatrix}A & 0\\0 & A\end{pmatrix},
\quad
D = A^2 > 0,\quad
\frac{\partial^2 E}{\partial\theta^2} = A > 0.
\]

Es un \textbf{mínimo local}: el panel orientado en dirección opuesta al Sol.

\subsubsection{Puntos de silla}

En $(\theta_0,\phi_0+\pi)$ y $(\theta_0+\pi,\phi_0)$:

\[
H_E(\theta_0,\phi_0+\pi)=\begin{pmatrix}-A & 0\\0 & A\end{pmatrix},
\quad
H_E(\theta_0+\pi,\phi_0)=\begin{pmatrix}A & 0\\0 & -A\end{pmatrix},
\quad
D = -A^2 < 0,
\]

ambos son \textbf{puntos de silla}.

\subsection{Curvas de nivel}

Las curvas de nivel son los conjuntos donde $E(\theta,\phi)=c$:

\begin{equation}\label{eq:curvas}
A\cos(\theta-\theta_0)\cos(\phi-\phi_0)=c
\quad\Longrightarrow\quad
\cos(\theta-\theta_0)\cos(\phi-\phi_0)=\frac{c}{A},
\end{equation}

con $|c|\leq A$ para que existan soluciones reales.

Interpretación:
\begin{itemize}
    \item $c=A$: solo el punto $(\theta_0,\phi_0)$ (máximo).
    \item $c$ cercano a $A$: regiones aproximadamente elípticas alrededor del óptimo, representando configuraciones de alto rendimiento.
    \item $c=0$: rectas $\theta-\theta_0=\pi/2$ y $\phi-\phi_0=\pi/2$, formando una cuadrícula que separa regiones de energía positiva y negativa.
\end{itemize}

\subsection{Derivada direccional}

La derivada direccional en la dirección del vector unitario $\mathbf{u}=(u_1,u_2)$ es

\begin{equation}\label{eq:dd}
D_{\mathbf{u}}E(\theta,\phi) = \nabla E(\theta,\phi)\cdot\mathbf{u}
= -A\sin(\theta-\theta_0)\cos(\phi-\phi_0)\,u_1
  -A\cos(\theta-\theta_0)\sin(\phi-\phi_0)\,u_2.
\end{equation}

Para $\mathbf{u}=(\cos\alpha,\sin\alpha)$, representa la variación energética al modificar simultáneamente inclinación y orientación según un ángulo $\alpha$ respecto al eje $\theta$.

\subsection{Plano tangente y linealización}

La linealización de $E$ en torno a $(\theta^*,\phi^*)$ es

\[
L(\theta,\phi)=E(\theta^*,\phi^*)
+\frac{\partial E}{\partial\theta}(\theta^*,\phi^*)\,(\theta-\theta^*)
+\frac{\partial E}{\partial\phi}(\theta^*,\phi^*)\,(\phi-\phi^*).
\]

En el óptimo $(\theta_0,\phi_0)$:

\[
L(\theta,\phi)=A + 0\cdot(\theta-\theta_0) + 0\cdot(\phi-\phi_0)=A.
\]

La aproximación de primer orden es constante alrededor del máximo, pues las derivadas parciales se anulan. Para capturar la curvatura se requiere la expansión de Taylor de segundo orden:

\[
E(\theta,\phi)\approx A
+\frac12\frac{\partial^2 E}{\partial\theta^2}(\theta_0,\phi_0)\,(\theta-\theta_0)^2
+\frac12\frac{\partial^2 E}{\partial\phi^2}(\theta_0,\phi_0)\,(\phi-\phi_0)^2
+\frac{\partial^2 E}{\partial\theta\partial\phi}(\theta_0,\phi_0)\,(\theta-\theta_0)(\phi-\phi_0).
\]

Evaluando las derivadas segundas en $(\theta_0,\phi_0)$:

\begin{equation}\label{eq:taylor2}
E(\theta,\phi)\approx A-\frac{A}{2}\big[(\theta-\theta_0)^2+(\phi-\phi_0)^2\big],
\end{equation}

válida para puntos cercanos al óptimo. La energía decrece aproximadamente como un paraboloide en ambas direcciones.

\subsection{Análisis de sensibilidad}

La sensibilidad del sistema ante errores de instalación se estudia mediante las derivadas parciales. Para pequeñas desviaciones $\Delta\theta$, $\Delta\phi$ respecto de la configuración óptima, la variación de energía es

\[
\Delta E \approx \left|\frac{\partial E}{\partial\theta}\right|\Delta\theta
+ \left|\frac{\partial E}{\partial\phi}\right|\Delta\phi.
\]

Cerca del óptimo las derivadas parciales se anulan, por lo que el término dominante proviene de la expansión de segundo orden \eqref{eq:taylor2}. Un error combinado de magnitud $\delta$ (es decir, $\sqrt{(\Delta\theta)^2+(\Delta\phi)^2}=\delta$) produce aproximadamente

\begin{equation}\label{eq:sens}
\Delta E \approx \frac{A}{2}\,\delta^2.
\end{equation}

Esto indica que el sistema es cuadráticamente tolerante a pequeños errores: un error de $5^\circ$ ($\delta\approx0.0873$ rad) produce una pérdida de solo $0.38\%$ de $A$.

%% ====================================================================
\section{Implementación computacional}
%% ====================================================================

\noindent\textit{Esta sección se completará una vez desarrollado el código computacional.}

La implementación contempla:
\begin{enumerate}
    \item Interfaz interactiva para modificar $\theta$ y $\phi$ y visualizar la energía captada.
    \item Cálculo de derivadas parciales $\partial E/\partial\theta$, $\partial E/\partial\phi$ y gradiente $\nabla E$.
    \item Gráfica de curvas de nivel $E(\theta,\phi)=c$.
    \item Gráfica de la superficie $z=E(\theta,\phi)$ con puntos críticos identificados.
    \item Optimización numérica y verificación de la solución analítica.
    \item Análisis de sensibilidad ante errores de instalación.
    \item Visualización del modelo extendido con factores adicionales.
\end{enumerate}

Se utilizará Python con \texttt{numpy}, \texttt{matplotlib} y \texttt{scipy}. El código fuente se entrega por separado.

%% ====================================================================
\section{Conclusiones}
%% ====================================================================

\noindent\textit{Esta secci\'on se completar\'a una vez obtenidos los resultados num\'ericos de la implementaci\'on computacional.}

Aspectos a abordar en la versión final:
\begin{enumerate}
    \item Verificación de la configuración óptima obtenida analíticamente.
    \item Influencia de la inclinación $\theta$ y la orientación $\phi$ sobre la energía.
    \item Comportamiento del gradiente en distintas regiones del dominio.
    \item Interpretación de las curvas de nivel obtenidas.
    \item Sensibilidad frente a errores de instalación.
    \item Comparación entre configuraciones óptimas y no óptimas.
    \item Factibilidad práctica de las soluciones encontradas.
    \item Cumplimiento de los objetivos específicos planteados.
    \item Limitaciones del modelo y posibles mejoras.
\end{enumerate}

%% ====================================================================
\bibliographystyle{plain}
\bibliography{refs}

\end{document}
